% ====================================================================
%  递交前注意事项（内部工作文档，非递交件）
%  由 notes.sty 排版。要改内容直接改本文件，这是唯一来源。
% ====================================================================
\documentclass[11pt,UTF8]{ctexart}
\usepackage{notes}
\notestitle{递交前注意事项}

\begin{document}

\nH{递交前注意事项与代理人核查清单（A1）}

\nHH{一、最高优先级：申请人自身在先公开}

代码仓库已于 {\ntbf 2026-09-13} 公开推送（\nc{origin/0.4.69}，提交 \nc{fedd467f}）。其中下列文件完整披露了本申请的技术方案：

\begin{itemize}
  \item \nc{crates/runtime/src/evidence\_ledger.rs}（双指纹、易变字段剔除、三重合取阻断、轮询倍增、从消息重建）
  \item \nc{crates/runtime/src/tool\_output\_artifact.rs}（外置阈值、内容摘要值命名、引用式结果、预览截取）
  \item \nc{crates/runtime/src/conversation.rs}（自适应外置阈值、拒绝结果不入账、指示注入）
  \item \nc{docs/development-logic/evidence-ledger.md}（方案的自然语言完整描述）
  \item \nc{CHANGELOG.md} v0.4.69 条目（含全部阈值常量的具体数值）
\end{itemize}

《专利法》第 24 条的六个月宽限期仅适用于：国家紧急状态或非常情况下为公共利益首次公开、中国政府主办或承认的国际展览会首次展出、规定的学术会议或技术会议首次发表、他人未经申请人同意泄露。{\ntbf 申请人自行在代码托管平台公开不属于上述任一情形，不享有宽限期。}

{\ntbf 后果}：自 2026-09-13 起，上述公开即构成本申请的现有技术，足以单独破坏新颖性。

{\ntbf 唯一缓解}：以最快速度递交。每推迟一日，风险只增不减。

{\ntbf 不建议}的做法：删除或私有化仓库。公开事实已经发生（含第三方镜像、搜索引擎缓存与克隆），事后删除既不能消灭现有技术，也可能在后续程序中造成不利印象。

\begin{nquote}
若申请人希望保留完整的优先权布局（PCT / US / EP），须知美国有 35 U.S.C. §102(b)(1) 的 12 个月自我公开宽限期，欧洲专利局则与中国类似、{\ntbf 不提供}一般性宽限期。因此 EP 路线大概率已经关闭，US 路线仍开放至 2027-09-13。
\end{nquote}

\nHH{二、次高风险：抵触申请（检索不可见）}

专利申请有 18 个月公开延迟。以 OpenAI / Anthropic / 微软 / 字节等主体为申请人的在先申请若在本申请日之前提出、在本申请日之后公开，构成《专利法》第 22 条第 2 款的抵触申请，破坏新颖性且任何检索都查不到。agent runtime 的循环防护是上述主体必然遇到的问题。该风险{\ntbf 无法规避，只能靠早申请}。

\nHH{三、说明书刻意未写入的内容及其原因}

\begin{ntable}{Y{0.46}Y{0.46}}
{\ntbf 未写入的内容} & {\ntbf 原因} \\
\nthead
具体的现有技术文献号（US20110280145A1、WO2025221751A1 等） & 本工具未独立核实这些文献号与其记载内容的对应关系。在申请文件中主张未经核实的文献构成事实性错误风险。文献清单见 \nc{patent/PRIOR\_ART.md}，{\ntbf 请代理人核实后再决定是否写入背景技术}。CN 实践中背景技术不引用文献号完全可以接受。 \\
任何实验数据、命中率、节省比例、性能对比 & 《专利审查指南》与本仓 \nc{patent-writing-principles} 的明确要求；此类内容属于论文而非专利。 \\
仓库名、模块路径、环境变量以外的实现标识 & 避免把保护范围绑定到某一具体实现。 \\
对现有技术优越性的承认 & 不得出现"现有技术工作良好，但……"式表述。 \\
\end{ntable}

\nHH{四、创造性论证的三个落点（答复审查意见时使用）}

审查员最可能的驳回路径是"把已知技术应用到新领域属于常规应用"（引证内容哈希 + 计数判环的网络包环路检测、爬虫去重、邮件去重剔除时间戳）。预备的答辩支点按强度排序：

{\ntbf 支点一（最强）——账本可从消息历史重建。}

现有的内容哈希判环技术运行于状态不会被销毁的环境（路由器转发表、爬虫队列、邮件服务器）。大语言模型智能体的独特技术前提是：承载判重状态的会话上下文本身会因超出预算而被压缩，状态随之丢失。"判重状态必须能够从已存储的调用—结果配对按同一规则完整重建"是该前提下才会出现的技术问题，现有技术既未提出该问题也未给出解法。{\ntbf 此为本申请最好的创造性支点，应在答辩中置于首位。}

{\ntbf 支点二——内容摘要值充当证据身份（外置与判重的协同）。}

本领域的上下文预算管理（MemGPT、WO2025221751A1）与循环检测是两条独立的技术线。二者简单并置时，外置会使判重失效：产物路径含调用标识，同一内容必落在不同路径上，判重把每个副本都当作新证据。本申请以内容摘要值而非路径作为证据身份，使二者可同时成立。这是{\ntbf 协同效应而非功能拼接}，符合创造性判断中"技术特征在功能上彼此支持"的要求。

{\ntbf 支点三——三重合取阻断条件。}

现有技术（含 US20110280145A1）均为"单一条件计数达 k 次"。本申请要求无新证据连续计数、调用指纹计数、同结果计数三者同时达阈值，并对轮询类工具放大连续计数阈值。其解决的是单条件方案在智能体场景下同时存在的误杀（正当轮询）与漏判（入参微变的同义调用）两个具体缺陷。

{\ntbf 技术效果的表述口径}：统一表述为"减少无效的工具调用与外部网络请求，降低算力与令牌消耗"，即《专利审查指南》认可的"提升计算机系统内部性能"。{\ntbf 不要}表述为"提高回答质量""让模型更聪明"——后者会把方案推向"智力活动的规则和方法"。

\nHH{五、权利要求核查清单（代理人）}

\begin{itemize}[label=\nbox]
  \item 权 1 的前序部分是否足够短？目前前序仅描述"模型—工具迭代循环"这一公知框架，其余全部落在特征部分。确认无发明点被误置于前序。
  \item 权 1 步骤 S3 已写为{\ntbf 肯定的方法步骤}（"对……超过外置阈值的工具调用，将该工具输出存储为产物文件"）而非条件从句，以避免被认定为可选特征。确认该写法与实施例四一致。
  \item 第一至第四阈值在权 1 中均以序数词给出、未写入具体数值，具体数值（4 / 6 / 3 / 3）仅出现在具体实施方式。确认此分层无误——数值写入权 1 会不必要地收窄范围。
  \item 权 1 未出现"其中"以外的多重嵌套定语，单句长度是否仍在审查员可接受范围内？若认为过长，可将 S2 的易变字段列举移入权 2。
  \item 权 13（系统）的各单元是否与说明书图 1 的 101～109 一一对应？目前 101～106、108 已在权 13 中出现，107（会话上下文存储单元）与 109（事件记录单元）未出现在权 13 中 —— 确认这是有意的（二者分别由权 1 的上下文描述与权 12 覆盖），或补入权 13。
  \item 权 14 引用"权利要求 1 至 12 中任一项"，未引用权 13（系统）。确认符合介质权项的常规写法。
  \item claims 连续编号 1–14，无断号，独立权与从属权交错编号（权 13、权 14 为独立权）。
  \item 超出 10 项的部分（权 11–14 共 4 项）产生附加费 4 × 150 = 600 元。确认申请人接受。
\end{itemize}

\nHH{六、说明书核查清单（代理人）}

\begin{itemize}[label=\nbox]
  \item 支持性：权 2～权 12 的每一个附加技术特征是否都能在具体实施方式中找到对应段落？
\end{itemize}

对照表——权 2→实施例二 S205 情形一；权 3→实施例三 S307；权 4→实施例二末段；权 5→实施例一末段；权 6→实施例四 S401；权 7→实施例四 S406；权 8→实施例四 S405、S404；权 9→实施例五 S501；权 10→实施例五 S505；权 11→实施例三末段；权 12→实施例二 S209。

\begin{itemize}[label=\nbox]
  \item 充分公开：本领域技术人员能否仅凭说明书实现？重点检查实施例二 S205（错误签名的生成）与实施例四 S406（预览的截取规则）是否已足够具体。
  \item 术语一致性：全文是否严格使用"调用指纹 / 证据指纹 / 无新证据连续计数 / 内容摘要值 / 产物文件 / 引用式结果 / 策略变更指示 / 轮询类工具"，未出现同义换词。术语表见 \nc{patent/INVENTION\_DISCLOSURE.md} 第六节。
  \item "所述"用法：首次引入用"一个/一种"，其后一律"所述"。
  \item 附图标记：说明书出现的 100～111 与 \nc{附图绘制说明.md} 是否完全一致，且图中无说明书未提及的标记。
  \item 摘要字数：已核，{\ntbf 297 字}（≤300，含标点）。摘要附图指定为图 1。
\end{itemize}

\nHH{七、本工具未执行的步骤}

\begin{itemize}
  \item {\ntbf 未执行外部审查员交叉复核}：\nc{codex} MCP 服务在本次会话中连接失败（\nc{MCP error -32000: Connection closed}），\nc{/patent-review} 阶段的 GPT-5.6-Sol 审查未运行。建议在递交前由代理人或重新连接 codex 后补做。
  \item {\ntbf 未进行专业检索}：\nc{patent/PRIOR\_ART.md} 中的检索结果由申请人提供，本工具未独立复核任何文献号。递交前应委托专业检索机构或代理人复核。
  \item {\ntbf 附图已生成，但未经制图人员复核}：\nc{说明书附图.html}（矢量 SVG）与 \nc{pdf/说明书附图.pdf}（5 页）为本工具按 \nc{附图绘制说明.md} 的规格绘制，黑白线条、无着色无阴影、除必需词语外无注释。递交前仍建议由代理人或制图人员按该文件的复核要点过一遍，尤其是图 1 中 107→101 与 108→104 的唯一一处连线交叉是否符合受理惯例。
\end{itemize}

\nHH{八、建议的递交顺序}

\begin{enumerate}
  \item 代理人复核权利要求书与说明书（半日至一日）；
  \item 制图人员按 \nc{附图绘制说明.md} 出 5 幅正式附图（一日）；
  \item 同日递交 CN 发明专利申请，取得申请日；
  \item 申请日取得后，再从容处理：是否同步递交实用新型（{\ntbf 不可行}——本申请为方法，实用新型不接受方法权利要求；若要覆盖，须另行撰写纯装置权项的实用新型，但其保护范围将显著窄于权 13）、是否在 12 个月内走 PCT 或 US 路线。
\end{enumerate}

\end{document}
